\documentclass[]{fairmeta}
% Option "twocolumn" available, but please prioritize single-column

% \usepackage{xcolor}
% \usepackage{fancyhdr,graphicx}
\usepackage{amssymb,amsthm,amsmath}
% \usepackage{booktabs,array} % for tables
% \usepackage{thmtools}
% \usepackage{thm-restate}
% \usepackage{algorithm,algorithmicx,algpseudocode} skip for now
% \usepackage[title]{appendix}
% \usepackage{tikz}
% \usepackage{resizegather}
% \usepackage{wrapfig}
% \usepackage{enumitem,kantlipsum}
\usepackage{mathtools}
% \usepackage{nccmath}
\usepackage{graphicx}



% \newtheoremstyle{mystyle}
% {6pt} % Space above
% {\topsep} % Space below
% {} % Body font
% {} % Indent amount
% {\bfseries} % Theorem head font
% {.} % Punctuation after theorem head
% {.5em} % Space after theorem head
% {} % Theorem head spec (can be left empty, meaning `normal')

% \theoremstyle{definition}
% \newtheorem{definition}{Definition}
% \newtheorem{assumption}{Modelling assumption}
% \theoremstyle{plain}
% \newtheorem{proposition}{Proposition}
% \newtheorem{example}{Example}
% \newtheorem{theorem}{Theorem}
% \newtheorem{corollary}{Corollary}
% \newtheorem{lemma}{Lemma}
% \newtheorem{observation}{Observation}

\input{macros}

%\usepackage[numbers,sort&compress]{natbib}
\usepackage{multicol}
\usepackage{enumitem}
% Text + math both in Times style
\usepackage{newtxtext}
\usepackage{newtxmath}

\usepackage{hyperref}
\usepackage{url}

\usepackage{float}

\usepackage[most]{tcolorbox}
\usepackage{xcolor}

\title{Religious Bias in LLMs is Significantly Understudied}

\author[1,\dagger]{Walter Reade}
\author[1,2]{Sheryl Carty}
\author[1,2]{Nancy Fulda}

\contribution[\dagger]{Corresponding author, inversion@kaggle.com}
\affiliation[1]{The Consortium for Evaluation of Faith and Ethics in AI}
\affiliation[2]{Brigham Young University}

\abstract{
In the earlier years of development of LLMs, it was relatively easy to prompt an LLM to respond with toxic or biased statements about religion. Subsequent improvements in frontier models addressed many of the issues of bias and toxicity in general, including against religion. At the same time, the adoption and usage of these models has grown exponentially. Small and implicit biases, therefore, have a magnified overall impact. In this paper, we (1) briefly review previous efforts to measure religious bias in LLMs, (2) show, by reviewing over 12,000 papers dealing with bias in LLMs, that religious bias has been significantly understudied compared to other bias targets, and (3) highlight areas where further research is needed to understand gaps in model performance for a growing global community of LLM users.
}

\date{\today}

\begin{document}

\maketitle

\section{Introduction}

\subsection{Motivation}

The rapid proliferation of Large Language Models into everyday life has created a new and unprecedented channel through which billions of people acquire information, seek advice, and form opinions. Early LLMs, unconstrained by safety mechanisms, generated openly racist content, offensive religious representations, and other forms of toxic output with alarming ease \citep{brown2020}. The introduction of safety measures, including Reinforcement Learning from Human Feedback (RLHF), constitutional AI, and content filtering, helped mitigate overt bias toward the groups most frequently targeted with hate speech \citep{ouyang2022, bai2022}.

However, the current landscape presents a subtler and potentially more consequential challenge. While existing biases are less explicit, the rapidly increasing number of individuals using LLMs for knowledge acquisition, decision-making, and even spiritual guidance significantly amplifies the impact of even subtle biases. In a world where approximately 3 out of 4 people align with a religious tradition \citep{pew}, a model that slightly favors one religious perspective over another, or that treats one tradition with nuance while reducing another to stereotypes, can shape the understanding of millions of users daily. The scale of deployment transforms what might be a minor imperfection in a research prototype into a systemic force with real-world consequences for religious communities worldwide.

\subsection{Objectives}

This paper pursues three primary objectives:
\begin{enumerate}
    \item Review previous studies measuring religious bias in LLMs, including a comprehensive analytical analysis of studies published to arXiv.org.
    \item Identify research gaps in how religious bias has been studied relative to other forms of demographic bias.
    \item Propose highest-impact next steps for future research.
\end{enumerate}

We take a pragmatic view in our recommendations. While there is a very long tail of possible research questions, we focus on those that have the most significant real-world impact on the greatest number of people. We prioritize the models people are actually using in their daily experience over models that may be more convenient to study. And we aim and encourage others to develop methodologies that can be extended by other religious communities to address the long tail of research opportunities.

\subsection{Brief Overview of Progress}

Research on bias in language technologies has evolved through several distinct phases, each building on the insights and limitations of the previous era.

\paragraph{Pre-LLM foundations.} \citet{caliskan2017} demonstrated that word embeddings trained on human-generated corpora inherit human-like biases, including associations between religious groups and stereotypical attributes. This foundational work established that language models are not neutral artifacts but cultural mirrors that reflect, and can amplify, the prejudices embedded in their training data.

\paragraph{Early LLM studies.} As generative language models grew in scale and capability, researchers began examining their outputs for bias. \citet{abid2021} conducted a landmark prompt-continuation study showing that when GPT-3 was given the prompt ``Two Muslims walked into a...,'' the model persistently generated completions associating Muslims with violence---completing the prompt with references to shootings, bombings, and other acts of terrorism in 66\% of cases, compared to only 15\% for analogous prompts about Christians. This work remains the most-cited paper focused primarily on religious bias in LLMs.

\paragraph{Standardized benchmarks.} To enable systematic evaluation, the research community developed standardized benchmarks including StereoSet \citep{nadeem2021}, CrowS-Pairs \citep{nangia2020}, BOLD \citep{dhamala2021}, and BBQ \citep{parrish2022}. These benchmarks operationalized bias measurement through techniques such as pseudo-log-likelihood scoring, counterfactual token swapping, and open-ended generation analysis. Each includes a religion category, enabling comparison across demographic dimensions. But, as we will show, religion has typically been included as one category among many rather than studied in depth.

\paragraph{Safety alignment and its side effects.} The introduction of RLHF and related techniques as mechanisms to reduce bias and toxicity in model outputs brought significant improvements but also introduced unintended consequences. Models began exhibiting what researchers have termed ``vacuous neutrality'', i.e., refusing to answer benign questions about religious topics or defaulting to uninformative ``Cannot be determined'' responses whenever religion is mentioned. \citet{santurkar2023} demonstrated that aligned LLMs disproportionately reflect the values of WEIRD (Western, Educated, Industrialized, Rich, Democratic) populations, a consequence of relying primarily on Western annotators for RLHF. This WEIRD bias has particular implications for religion, where the world's diversity of belief systems, practices, and theological frameworks far exceeds the range of perspectives represented in alignment training data.


\section{Analysis of arXiv: Methodology}
\label{sec:Methodology}

\subsection{Overview}

To quantify the state of religious bias research in LLMs, we conducted a comprehensive analysis of the arXiv preprint repository. We chose arXiv for the scope of this analysis because of its breadth (including, e.g., coverage it provided to conference papers), because it generally contains the most up-to-date studies, and because the scope allows for meaningful statistical comparison.

To identify qualifying papers, rather than relying on keyword-based searches, which we found to be unreliable due to inconsistent terminology across studies, we employed a multi-stage AI-assisted pipeline to systematically identify, classify, and analyze relevant research.

\subsection{Pipeline Design}

Our analysis proceeded through four stages:

\paragraph{Stage 1: Broad identification.} Every abstract on arXiv.org was assessed by Gemini Flash to identify studies dealing with bias, fairness, ethics, or morality in NLP and large language models. This comprehensive, brute-force approach was necessitated by the unreliability of keyword searches: papers studying religious bias use widely varying terminology (e.g., ``faith,'' ``denomination,'' ``sectarian,'' ``Islamophobia,'' ``antisemitism''), and many relevant papers do not include the word ``religion'' or ``bias'' in their abstracts despite substantively measuring religious bias.

\paragraph{Stage 2: Full-paper metadata extraction.} For each paper identified in Stage 1, Gemini Flash processed the complete contents of the paper to extract metadata (page count, number of references, publication date) and, critically, to determine whether the paper specifically measured bias and fairness in LLMs, as opposed to merely discussing the topic.

\paragraph{Stage 3: Bias target extraction.} Papers confirmed as measuring bias and fairness in LLMs were analyzed using Gemini Pro to extract the specific list of bias targets measured (e.g., gender, racial/ethnic, religious, age) and whether one of these targets was the primary focus of the study.

\paragraph{Stage 4: Deep religious bias analysis.} Papers identified as measuring religious bias were subjected to a detailed analysis using Gemini Pro to extract:
\begin{itemize}
    \item \textbf{Measurement description:} What specifically the paper measured or evaluated in terms of faith and religion.
    \item \textbf{Religious groups studied:} Which specific religious groups were measured or mentioned (e.g., Christianity, Islam, Judaism, Buddhism, Hinduism, Atheism).
    \item \textbf{Models tested:} Which specific LLMs were evaluated.
    \item \textbf{Languages evaluated:} Which languages were used in the evaluation (English assumed if not stated).
    \item \textbf{Religion component:} The extent to which religion is a focus of the paper (``major'' if religion is the paper's main focus, ``minor'' if it is a component but not the main focus).
    \item \textbf{Base benchmarks:} Existing benchmarks the paper extends or compares against.
    \item \textbf{Key findings:} The primary results related to religion.
    \item \textbf{References:} All references cited in the paper.
\end{itemize}


\section{Analysis of arXiv: Results}

\subsection{The Landscape of Bias Research in LLMs}

 Table~\ref{tab:paper_counts} shows the counts of papers on arXiv matching the specific categories:

\begin{table}[htbp]
\centering
\caption{Paper counts extracted from arXiv using the methodology in Section \ref{sec:Methodology}. Out of 12,851 papers identified as addressing bias and fairness in LLMs, only 34 (0.26\%) engage deeply with the topic of religious bias.}
\label{tab:paper_counts}
\begin{tabular}{lr}
\toprule
\textbf{Category} & \textbf{Count} \\
\midrule
Papers on arXiv (total corpus) & 3,043,207 \\
Papers about LLMs & 167,291  \\
Papers measuring bias and fairness in LLMs & 12,851 \\
Papers measuring religious bias & 1,374 \\
Papers where religious bias is the primary focus & 34 \\
\bottomrule
\end{tabular}
\end{table}

The funnel from papers measuring any form of bias in LLMs (12,851) to those measuring religious bias (1,374, or 10.7\%) to those where religion is the primary focus (34, or 0.26\%) reveals a striking attenuation. While many papers have measured religious bias, it has been measured almost exclusively as part of broader sweeps across demographic categories using existing benchmark datasets (e.g., StereoSet, CrowS-Pairs, BBQ), and rarely as the main focus of dedicated study.

\subsection{Religious Bias in Context: Comparison Across Bias Categories}

To contextualize the relative under-representation of religious bias research, we compare the frequency with which different bias targets appear across the 12,851 papers, both as any mention and as a primary focus, shown in Table~\ref{tab:bias_targets}.

\begin{table}[htbp]
\centering
\caption{Top bias targets measured across 12,851 LLM bias papers. Religious bias is often included as a sub-component in bias benchmarks, but is only rarely the subject of nuanced and systematic study.}
\label{tab:bias_targets}
\begin{tabular}{lrrrr}
\toprule
\textbf{Bias Target} & \textbf{Papers Measuring} & \textbf{\% of Total} & \textbf{Primary Focus} & \textbf{\% of Total} \\
\midrule
Gender bias & 3,927 & 30.8\% & 902 & 7.1\% \\
Racial/Ethnic bias & 2,612 & 20.5\% & 145 & 1.1\% \\
Language bias & 1,677 & 13.2\% & 760 & 6.0\% \\
Cultural bias & 1,364 & 10.7\% & 600 & 4.7\% \\
Religious bias & 1,336 & 10.5\% & 37 & 0.3\% \\
Age bias & 1,324 & 10.4\% & --- & --- \\
Political bias & 1,162 & 9.1\% & 376 & 2.9\% \\
Sexual orientation bias & 906 & 7.1\% & --- & --- \\
Socioeconomic bias & 871 & 6.8\% & 41 & 0.3\% \\
Disability bias & 766 & 6.0\% & 80 & 0.6\% \\
\bottomrule
\end{tabular}
\end{table}

Religious bias is the 5th most frequently measured category overall (10.5\%), appearing at roughly the same rate as cultural and age bias. Yet when we examine primary focus studies, religion drops dramatically. Only 37 papers (0.26\% of the total) primarily focus on religious bias, the lowest ratio of primary-focus to incidental-measurement of any major bias category. For comparison, gender bias is the primary focus of 23\% of the papers that measure it, while religious bias is the primary focus of only 2.8\% of the papers that measure it.

This finding is the central quantitative result of our analysis: religious bias is routinely measured but rarely studied in depth.

\subsection{Religious Groups Studied}

Among the 1,374 papers that measure religious bias, the distribution of religious groups studied reveals significant imbalances, as shown in Table~\ref{tab:religious_groups}.

\begin{table}[htbp]
\centering
\caption{Religious groups studied across 1,374 papers. Abrahamic religions receive the most attention. Minority religions, named denominations or major religions, and various global religions are far less represented.}
\label{tab:religious_groups}
\begin{tabular}{lrr}
\toprule
\textbf{Religious Group} & \textbf{Count} & \textbf{\% of Papers} \\
\midrule
Islam & 882 & 64.2\% \\
Christianity & 598 & 43.5\% \\
Judaism & 576 & 41.9\% \\
Hinduism & 354 & 25.8\% \\
Buddhism & 315 & 22.9\% \\
Atheism & 212 & 15.4\% \\
Religion (general) & 159 & 11.6\% \\
Catholic & 108 & 7.9\% \\
Sikhism & 107 & 7.8\% \\
Protestant & 64 & 4.7\% \\
Mormon/Latter-day Saint & 60 & 4.4\% \\
Agnosticism & 43 & 3.1\% \\
Taoism & 38 & 2.8\% \\
Jainism & 28 & 2.0\% \\
Shinto & 18 & 1.3\% \\
Sunni & 17 & 1.2\% \\
Zoroastrianism & 14 & 1.0\% \\
Shia & 14 & 1.0\% \\
Paganism & 12 & 0.9\% \\
Confucianism & 12 & 0.9\% \\
\bottomrule
\end{tabular}
\end{table}

\paragraph{Abrahamic dominance.} The vast majority of studies include the three Abrahamic religions: Islam (64.2\%), Christianity (43.5\%), and Judaism (41.9\%). Islam's prominence as the most-studied group reflects the severity and persistence of anti-Muslim bias documented across language models \citep{abid2021, naous2024}.

\paragraph{Denominational erasure.} Specific denominations are studied significantly less frequently than umbrella classifications. While ``Christianity'' appears in 43.5\% of papers, ``Catholic'' appears in only 7.9\% and ``Protestant'' in 4.7\%. The same pattern holds within Islam: while the umbrella term appears in 64.2\% of papers, ``Sunni'' and ``Shia'' appear in only 1.2\% and 1.0\%, respectively. This monolithic treatment obscures the vast theological, cultural, and practical differences between traditions within each religion. As one recent study \citep{maheshwari2026} demonstrated, frontier models exhibit significant intra-religious biases; for example, DeepSeek-v3 and GPT-4o favored Shia interpretations in English but switched to favoring Sunni viewpoints when the same questions were posed in Hindi.

\paragraph{Minority religions are effectively invisible.} Religions with hundreds of millions of adherents globally---including Sikhism, Taoism, and Jainism---receive minimal attention (7.8\%, 2.8\%, and 2.0\% respectively). Groups such as Zoroastrianism, Paganism, and Confucianism each appear in approximately 1\% of papers or fewer.

\subsection{Languages Evaluated}

The linguistic landscape of religious bias research is overwhelmingly monolingual, as shown in Table~\ref{tab:languages}.

\begin{table}[htbp]
\centering
\caption{Languages evaluated across 1,374 studies. Out of 12,851 LLM bias papers extracted from arXiv, less than 7\% explore LLM behavior in languages other than English.}
\label{tab:languages}
\begin{tabular}{lrr}
\toprule
\textbf{Language} & \textbf{Count} & \textbf{\% of Papers} \\
\midrule
English & 1,282 & 93.3\% \\
Chinese & 123 & 9.0\% \\
Spanish & 84 & 6.1\% \\
German & 80 & 5.8\% \\
French & 77 & 5.6\% \\
Arabic & 71 & 5.2\% \\
Korean & 58 & 4.2\% \\
Hindi & 58 & 4.2\% \\
Bengali & 35 & 2.5\% \\
Indonesian & 30 & 2.2\% \\
\bottomrule
\end{tabular}
\end{table}

English dominates at 93.3\%, with the next most common language (Chinese) appearing in only 9.0\% of studies. Arabic, a language intimately connected to Islam, the most-studied religious group, is evaluated in only 5.2\% of papers. Hindi, the primary language of hundreds of millions of Hindus, appears in just 4.2\%. This English-centricity is particularly concerning because religious bias has been shown to shift depending on the language of prompting. For example, SectEval \citep{maheshwari2026} found that models changed their sectarian preferences between English and Hindi, and recent studies have shown that models prompted in Bengali exhibit systematic biases favoring Hindu-associated cultural terms over Muslim-associated ones in a predominantly Muslim-speaking population \cite{kabir2026xcrbench}.

\subsection{Models Tested}

The distribution of models evaluated in religious bias research reflects both the evolution of the field and the accessibility of different model families, as shown in Table~\ref{tab:models}.

\begin{table}[htbp]
\centering
\caption{Models most frequently tested (top 15). Notably absent from this list are models from LLM providers with the most global market share, including Google, Anthropic, Perplexity, xAI, and DeepSeek.}
\label{tab:models}
\begin{tabular}{lrr}
\toprule
\textbf{Model} & \textbf{Count} & \textbf{\% of Papers} \\
\midrule
GPT-4 & 179 & 13.0\% \\
GPT-4o & 171 & 12.4\% \\
BERT & 169 & 12.3\% \\
GPT-2 & 108 & 7.9\% \\
GPT-3.5 & 97 & 7.1\% \\
RoBERTa & 91 & 6.6\% \\
GPT-3 & 80 & 5.8\% \\
ChatGPT & 72 & 5.2\% \\
Llama-3.1-8B-Instruct & 45 & 3.3\% \\
Mistral-7B & 43 & 3.1\% \\
DistilBERT & 42 & 3.1\% \\
GPT-4o-mini & 41 & 3.0\% \\
Qwen2.5-7B-Instruct & 34 & 2.5\% \\
T5 & 33 & 2.4\% \\
GPT-3.5-Turbo & 31 & 2.3\% \\
\bottomrule
\end{tabular}
\end{table}

Two observations are noteworthy. First, a substantial portion of the literature (BERT at 12.3\%, GPT-2 at 7.9\%) evaluates models that are no longer representative of frontier capabilities, reflecting the time lag inherent in academic publishing. Second, there is a notable gap between the models most frequently studied and the models most widely used by consumers. While GPT-4 and GPT-4o are well-represented, models from Anthropic (Claude), Google (Gemini), and other providers that serve large user populations are underrepresented in the literature.

\subsection{Benchmarks Used}

The benchmarks most frequently employed to measure religious bias reveal a high degree of concentration, as shown in Table~\ref{tab:benchmarks}.

\begin{table}[htbp]
\centering
\caption{Most frequently used benchmarks. Out of 1,374 arXiv papers that examine religious bias, more than 37\% of papers rely on the same four benchmarks.}
\label{tab:benchmarks}
\begin{tabular}{lrr}
\toprule
\textbf{Benchmark} & \textbf{Count} & \textbf{\% of Papers} \\
\midrule
StereoSet & 207 & 15.1\% \\
CrowS-Pairs & 164 & 11.9\% \\
BBQ & 164 & 11.9\% \\
BOLD & 82 & 6.0\% \\
RealToxicityPrompts & 47 & 3.4\% \\
MMLU & 46 & 3.3\% \\
WinoBias & 46 & 3.3\% \\
TruthfulQA & 41 & 3.0\% \\
ToxiGen & 39 & 2.8\% \\
\bottomrule
\end{tabular}
\end{table}

Four benchmarks---StereoSet, CrowS-Pairs, BBQ, and BOLD---account for more than 37\% of religious bias evaluations. While these benchmarks have been instrumental in establishing the field, their repeated reuse raises concerns about benchmark saturation and the possibility that models have been optimized to perform well on these specific tests without achieving genuine bias reduction.


\section{Key Findings from the Literature}

Drawing on the corpus of 1,374 papers that measure religious bias, we synthesize the major findings into several thematic categories.

\subsection{Implicit Bias Outlives Explicit Toxicity}

While frontier models such as GPT-4 and Claude 3 will refuse prompts like ``Write a hateful tweet about Jews,'' they continue to fail implicit bias tests. When presented with ambiguous scenarios, e.g., a crime committed by either a Christian or a Muslim, models disproportionately use Chain-of-Thought reasoning to justify attributing blame to the marginalized religious group based on latent stereotypes \cite{wu2025reasoning,turpin2023language}. Even after explicit debiasing, the association between Muslims and violence persists. \citet{abid2021} found this pattern in GPT-3, and subsequent work has confirmed its persistence: fine-tuning the Instruct series GPT-3 reduced explicit toxicity but left intact a ``second-order bias'' where common Muslim names in neutral prompts significantly increased violent completions compared to Christian or Hindu names ($p < .001$)---specifically, 18 violent completions for Muslim names versus 8 for Christian and 4 for Hindu names \cite{hemmatian2022debiased,hemmatian2023muslim}.

\subsection{The Safety Penalty and Over-Censorship}

A recurring finding across multiple papers is that safety-alignment mechanisms often harm minority religions disproportionately. Toxicity classifiers and LLM guardrails frequently rely on spurious correlations, associating identity terms like ``Muslim'' or ``Jewish'' with hate speech. This produces ``false refusals'', i.e., benign questions about minority religions are blocked because the safety filters misinterpret the mention of a religion as itself toxic \cite{chehbouni2024representational, plazadelarco2025false}. One study found that hate speech detection models achieve their highest accuracy (94.0\%) for content targeting Jewish people but drop to 75.3\% for content targeting Muslims, demonstrating that current safety guardrails create a hierarchy of protection that leaves certain communities more vulnerable than others \cite{selvaganapathy2026confident}.

This over-censorship has led to what researchers call ``vacuous neutrality'': Models trained to output ``Cannot be determined'' whenever religion is mentioned \cite{manduru2026beyond}. While this strategy improves benchmark scores, it results in a fundamental loss of utility and fails to align the model with genuinely pluralistic human values.

\subsection{Multimodal Vulnerabilities}

Bias in language models extends beyond text. Vision-Language Models (VLMs) consistently exoticize and stereotype religious groups---for example, depicting all Muslims in hijabs or skullcaps regardless of context, or generating Christian churches when prompted for a generic ``place of worship.'' \cite{ghosh2025documenting, qadri2023regimes}
 Text-to-image models are particularly vulnerable. \citet{zhang2026} introduced TwoHamsters, a benchmark demonstrating that individually benign words can be combined in prompts to bypass safety filters and create offensive religious imagery. Prompt combinations like ``Mosque + Pig'' or ``Crucifix + Trash Can'' produced offensive content at alarming rates: SD-XL achieved only a 6.90\% defense rate against the ``Mosque + Pig'' combination, meaning the model generated offensive imagery 93\% of the time. Current concept erasure methods fail to mitigate these risks, as the offensive meaning emerges from the composition of benign terms rather than from any single unsafe concept \cite{zhang2026}.

Furthermore, multimodal ``pragmatic jailbreaks'', where individually safe text and image elements combine to form harmful content, achieve attack success rates for religious hate speech of up to 72.0\% in some models, with an average success rate of 36.9\% across nine evaluated models \cite{liu2024multimodal}.

\subsection{WEIRD Bias and Cultural Defaults}

\citet{santurkar2023} demonstrated that aligned LLMs disproportionately reflect the values of WEIRD populations---a consequence of using primarily Western annotators for RLHF. This finding has profound implications for religious bias. Models default to Western, secular frameworks when discussing religion:
\begin{itemize}
    \item When evaluated against a Christian-specific understanding of human flourishing, all 20 evaluated frontier models defaulted to ``Procedural Secularism,'' prioritizing individual autonomy over theological reasoning, with an average 17-point performance decline and a 31-point decline in the Faith and Spirituality dimension specifically.
    \item PaLM's responses to World Values Survey items on religion leaned strongly toward US-centric religious positions rather than secular European ones.
    \item GPT-3 distorted the French concept of secularism (laïcité), reframing it through an American culture-war lens that conflated it with Islamophobia.
    \item A global survey found that ``Religion or Tradition'' was identified by 48\% of respondents as the primary anchor of cultural identity, with participants frequently establishing ``cultural redlines'' around sacred texts, prayers, and religious artifacts that they believe AI should never generate.
\end{itemize}

\subsection{Bias Amplification and Cross-Contamination}

Several studies have uncovered concerning dynamics in how bias propagates and interacts across categories:
\begin{itemize}
    \item \textbf{Chain-of-Thought amplification:} Standard Retrieval-Augmented Generation (RAG) generally reduces religious bias, but integrating Chain-of-Thought reasoning with RAG paradoxically amplifies it, increasing religion bias scores from 3.90 to 4.58--4.63, often exceeding the levels of the original model \cite{parihar2026evaluating}.
    \item \textbf{Single-edit catastrophic degradation:} Injecting just one biased sentence into an LLM via model editing (e.g., the ROME method) can cause catastrophic degradation across unrelated categories. A single religion-based stereotype injected into Llama3-8b increased its Gender Bias Score from 63.8\% to 81.6\% and its Race Bias Score from 29.6\% to 44.4\% \cite{chen2026editing}.
    \item \textbf{Persona-induced reasoning collapse:} Assigning a ``Religious'' persona to an LLM can cause up to a 69\% relative drop in accuracy on objective reasoning tasks like college chemistry, exposing deeply embedded stereotypes about the perceived scientific aptitude of religious individuals \cite{gupta2024bias}.
    \item \textbf{Quantization unpredictability:} Post-training quantization can unpredictably alter a model's treatment of specific religious demographics. After compressing LLaMA 3.2 3B, 15\% of responses concerning Catholics flipped their bias state, resulting in a statistically significant 8.4\% net increase in biased outputs \cite{hua2026uncertainty}.
\end{itemize}


\section{Research Gaps and Opportunities}

Our analysis reveals several critical areas where research is insufficient relative to the scale and importance of the problem.

\subsection{Non-Western and Multilingual Contexts}

With 93.3\% of religious bias studies conducted in English, the field has a massive blind spot for how models perform in the languages spoken by the majority of the world's religious practitioners. Our data shows that Arabic---the liturgical language of Islam, the most-studied religious group---is evaluated in only 5.2\% of papers. Hindi, spoken by hundreds of millions of Hindus, appears in just 4.2\%. Bengali, the language of over 230 million people in a region where Hindu-Muslim relations are a central societal concern, appears in only 2.5\%.

This gap is not merely academic. Recent research has demonstrated that religious bias shifts depending on the language of the prompt. Models exhibit different anti-Muslim biases when prompted in Hindi versus English; they favor Hindu-associated cultural terms over Muslim-associated ones in Bengali despite Muslims being the majority of Bengali speakers; and frontier models change their sectarian preferences between Sunni and Shia Islam depending on whether questions are asked in English, Hindi, or Arabic.

Furthermore, the question of localization extends beyond language to geographic context. How do models perform when queries are sourced from different geographic locations or IP addresses? Initial evidence from one study found that Claude Haiku completely flipped its loan approval preference based on spatiotemporal context, selecting a Lutheran-coded applicant 100\% of the time in ``California 1970'' but shifting to choose a Hindu-coded applicant 100\% of the time in ``Dearborn 2024.'' \cite{basu2026contextual}

\subsection{Intersectionality}

The intersection of religion with other demographic categories---gender, race, nationality, socioeconomic status---remains critically understudied. One paper found that inserting religious identity markers into clinical medical vignettes, both alone and intersectionally with sexual orientation, produced ``idiosyncratic, non-additive harms'' that systematically degraded model accuracy \cite{testoni2026calibrated}. Another found that religion serves as the most influential feature for predicting ethnicity in generated character profiles: removing the religion feature caused ethnicity classification accuracy to drop by up to 11.4\%, demonstrating models' overreliance on religious stereotypes as proxies for other demographic categories \cite{siddique2024who}.

This research is complicated by the combinatorial increase as groups are added, both increasing the cost of study and the complexity of analysis. Nevertheless, the real-world experience of religious individuals is inherently intersectional, and single-axis evaluations may systematically underestimate the biases they face.

\subsection{Multimodal Religious Bias}

Text-to-image safeguards are critical due to the ease with which they could be used to create offensive and hateful imagery. This area has received insufficient attention with regard to religious imagery, and the studies that have been performed demonstrate significantly weaker safeguards compared to text-generative models. \citet{zhang2026} demonstrated that two individually benign words could be combined in a prompt to bypass safety filters and create offensive images---``Mosque'' and ``Pig'' produced offensive content with defense rates as low as 3.45\%. Vision-Language Models consistently rendered Muslim Indians in stereotyped religious attire while depicting Hindu Indians in context-appropriate clothing for the exact same activities. Prompts for ``a house of worship'' generated Christian, American-looking churches, while prompts for ``Indian houses of worship'' lacked religious diversity, reflecting a ``Hinduization of Indian religious iconography.''

\subsection{Representational Norms: How Should Models Respond?}

We have found only two papers in our corpus that explicitly address the question of how LLMs should respond to religious queries, as opposed to how they do respond \cite{skytland2026evaluating, lutz2025goodfaith}. The conversation in the field has been dominated by bias detection and toxicity reduction, but the broader question of religious representation within LLM responses---the framing, depth, and perspective with which models engage religious content---is a much larger topic.

Silicon Valley appears to have adopted a de facto stance that religion should be avoided by default. Models are tuned to refuse engagement, deliver ``vacuously neutral'' responses, or frame religious topics through a secular lens \cite{skytland2026evaluating, lutz2025goodfaith}. We believe the research community should have an intentional conversation about whether this is the correct default behavior, particularly as LLMs become the primary information source for billions of people, many of whom hold religious belief as central to their identity.

\subsection{Ecological Validity of Evaluations}

Controlled prompts that are convenient for systematic analysis do not necessarily reflect how people actually interact with LLMs. The vast majority of evaluations use single-turn, isolated prompts (86.4\% of papers in our corpus do not employ continuous or multi-turn testing), yet real-world interactions are overwhelmingly multi-turn conversations where context accumulates and biases may compound. One study found that implicit biases injected into long-term memory ``accumulate and propagate to other unrelated social domains''---religion bias could increase biases in race or nationality over the course of extended interactions \cite{ma2026implicit}.

Furthermore, while stereotyping, implicit bias, and toxicity have been studied, there has been far less investigation into:
\begin{itemize}
    \item \textbf{Religious representation bias:} Whether models present religions with equal depth, nuance, and respect.
    \item \textbf{Moral and value misalignment:} Whether models accurately represent the ethical frameworks of different religious traditions rather than defaulting to secular Western ethics.
    \item \textbf{Reasoning quality:} The quality of model reasoning about religious topics, not merely the correctness of final answers. One study found that models exhibit ``sycophantic behavior by confidently accepting and justifying false religious premises'' with hallucination rates exceeding 55\% \cite{elmahjub2026islamiclegalbench}.
\end{itemize}

\subsection{Benchmark Integrity and Longitudinal Tracking}

The concentration of evaluations around a small number of benchmarks (StereoSet, CrowS-Pairs, BBQ, and BOLD collectively account for nearly half of all religious bias evaluations) raises critical concerns:
\begin{itemize}
    \item \textbf{Benchmark memorization:} Numerous benchmarks have been shown to have been memorized by LLMs, degrading their assessment fidelity \cite{white2025livebench}. Models may ``game'' multiple-choice fairness benchmarks by defaulting to refusal or neutral answers \cite{miandoab2025breaking}, masking latent biases that emerge in open-ended, real-world generation \cite{himelstein2025silenced}.
    \item \textbf{Need for hidden test sets:} To maintain evaluation integrity, the field needs hidden or dynamically generated test sets that cannot be incorporated into training data.
    \item \textbf{Longitudinal tracking:} Only 13.6\% of papers in our corpus employ any form of continuous testing. There is a critical need for more systematic longitudinal tracking across model versions to understand whether bias is genuinely decreasing or merely being displaced.
\end{itemize}

\subsection{Bias Interaction and Impossibility Results}

Some studies have shown that mitigating one bias can amplify another. A single-edit injection of a religion-based stereotype increased gender and race bias scores dramatically. This finding echoes the foundational work of \citet{kleinberg2016}, who proved that it is mathematically impossible to satisfy multiple commonly used definitions of fairness simultaneously in AI systems if base rates of an outcome differ between populations. Whether an analogous impossibility result applies to bias in LLMs---where reducing religious bias might necessarily increase bias along other dimensions---is an open and important question.

\subsection{User Trust and Reputational Harm}

To what degree do people accept AI-generated claims about religion as factual? This question is a critical component of understanding the real-world impact of religious bias. If users uncritically adopt LLM outputs as authoritative, then even minor biases in how a religion is represented can cause significant reputational harm to religious communities. The severity of this harm is amplified by the scale of deployment: a single model serving hundreds of millions of users can shape public perception more broadly than any individual media outlet.

\subsection{Additional Emerging Concerns}

Several additional research directions warrant attention:
\begin{itemize}
    \item \textbf{LLM-as-Judge frameworks:} As evaluation becomes more complex and the scope of assessment increases, better frameworks are needed for using LLMs to evaluate other LLMs' religious bias, with appropriate calibration and validation.
    \item \textbf{Group dynamics:} The introduction of group chat features in consumer LLM products (e.g., ChatGPT) and the integration of LLMs into social media platforms raises questions about group dynamics surrounding religious conversations. Do LLMs push back against stereotypes and bias expressed by humans? Do they support minority religious views in group contexts?
    \item \textbf{Situatedness:} How does the pattern of bias change when a model is provided additional context from chat histories, email exchanges, self-disclosed user information, or organizational data? How does it change when adopting a persona? One study found that religious personas introduced ``significant volatility'' in agent performance, with different religious identities improving or degrading performance depending on the model architecture \cite{cao2026biased}.
\end{itemize}


\section{Toward Comprehensive Evaluation: A Research Agenda}

Based on our analysis, we propose a research agenda organized by impact and feasibility:

\subsection{Immediate Priorities}

\begin{enumerate}
    \item \textbf{Dedicated religious bias benchmarks.} The field needs benchmarks designed specifically for religious bias evaluation---not merely religion categories within general-purpose bias benchmarks. These should cover the full spectrum of religious traditions, not just Abrahamic faiths, and should be developed in consultation with religious communities.
    \item \textbf{Multilingual evaluation.} Religious bias evaluation must be conducted in the languages spoken by religious practitioners, not just English. Priority languages should include Arabic, Hindi, Indonesian, Bengali, and Turkish, reflecting the linguistic diversity of the world's largest religious populations.
    \item \textbf{Frontier model coverage.} Evaluations must prioritize the models people actually use---including Claude, Gemini, and other widely deployed systems---rather than focusing disproportionately on models that are convenient for academic study.
\end{enumerate}

\subsection{Medium-Term Goals}

\begin{enumerate}[resume]
    \item \textbf{Multimodal safety audits.} Text-to-image and vision-language models require dedicated safety audits for religious content, with particular attention to compositional attacks (benign terms combining to produce offensive content).
    \item \textbf{Multi-turn and ecological evaluations.} Move beyond single-turn controlled prompts to evaluate bias in extended, naturalistic conversations that reflect how people actually interact with LLMs about religious topics.
    \item \textbf{Intersectional analysis.} Develop scalable methodologies for evaluating religious bias at the intersection of other demographic categories.
\end{enumerate}

\subsection{Long-Term Vision}

\begin{enumerate}[resume]
    \item \textbf{Representational norms.} Convene a multi-stakeholder conversation about how LLMs should engage with religious content---moving beyond avoidance toward respectful, accurate, and pluralistic engagement.
    \item \textbf{Community-extensible methodologies.} Design evaluation frameworks that can be extended by specific religious communities to address the long tail of traditions currently invisible in the literature.
    \item \textbf{Longitudinal monitoring.} Establish infrastructure for continuous, longitudinal tracking of religious bias across model versions, languages, and modalities.
\end{enumerate}


\section{Conclusion}

This paper has presented the most comprehensive analysis to date of religious bias research in Large Language Models. By systematically reviewing over 12,000 papers measuring bias and fairness in LLMs, we have demonstrated that religious bias occupies a paradoxical position: it is the 5th most frequently measured category of bias, yet it is the primary focus of only 0.3\% of papers, the lowest ratio of any major bias category. The 34 papers that primarily focus on religious bias comprises a tiny fraction of attention devoted to an issue that affects billions of religious individuals worldwide.

Our findings paint a picture of a research landscape that is heavily concentrated along multiple dimensions: on Abrahamic religions (with Islam, Christianity, and Judaism dominating), on English-language evaluation, on a small set of recycled benchmarks, and on text-only modalities. These concentrations leave critical gaps in our understanding of how LLMs treat the world's religious diversity.

The stakes are high. LLMs have evolved from research curiosities to infrastructure and are now embedded in search engines, education platforms, healthcare systems, and personal assistants used by billions. Small, implicit biases at this scale are not small in their effects. The persistent association of Muslims with violence, the over-censorship of minority religions, the defaulting to Western secular frameworks, and the vulnerability of multimodal systems to generating offensive religious imagery are not merely academic concerns. They are active harms affecting real communities.

We call on the research community, model developers, and policymakers to recognize religious bias as a first-class concern in AI safety---not merely an afterthought appended to gender and racial bias evaluations---and to invest in the dedicated, multilingual, multimodal, and community-engaged research that the scale of the problem demands.


\section*{Acknowledgments}
We thank the BH Roberts Foundation for their generous financial support of this project. We also gratefully acknowledge the contributions of CEFE\textperiodcentered{}AI research groups at Baylor University, Notre Dame University, and Yeshiva University.

\section*{About CEFE\textperiodcentered{}AI}

The Consortium for the Evaluation of Faith and Ethics in AI is a pluralistic, multi-university consortium of faith-based and research institutions committed to ensuring that religious representation in AI is honest, accurate, and respectful. Participating universities contribute faith-specific expertise, question design, and evaluation rubrics, while shared infrastructure supports standardized benchmarking across traditions. Consortium members include faculty from Baylor University, Brigham Young University, University of Notre Dame, and Yeshiva University. More information is available at \url{https://cefeai.org/}.

\appendix
\input{appendix}

\bibliographystyle{unsrtnat}
\bibliography{references}

\end{document}
